% AUTHORED
\section{Embeddings / token vectors}\label{sec:07}
\conceptbadge{07}

Neural networks operate on numbers, but language arrives as symbols. A tokenizer
maps text to integer token IDs. An embedding layer maps those IDs to trainable
vectors. If the vocabulary size is $V$ and the embedding dimension is $C$ (the same
channel dimension $C$ from \S00), the embedding table is
\[
  E \in \mathbb{R}^{V \times C}.
\]
For a token ID $i$, the embedding is simply the row $E_i$.

\subsection*{Lookup, not one-hot multiplication}
You can think of an embedding lookup as multiplying a one-hot vector by $E$, but
implementing it that way is wasteful. The one-hot vector contains one $1$ and
many zeros, so the result is just one row. In numpy, embeddings are indexing.
For token IDs with shape $B \times T$, the output has shape
$B \times T \times C$.

\subsection*{Training token vectors}
Embedding rows are parameters. During backpropagation, only the rows used by the
current batch receive nonzero gradient contributions. Frequent tokens are
updated often; rare tokens are updated only when they appear. Over training,
tokens that are useful in similar contexts tend to develop related vectors,
because the loss pushes them toward representations that support similar
predictions.

\subsection*{Character language model baseline}
The project starts with a character-level tokenizer. A simple baseline can use
the current character's embedding to predict the next character. This is not yet
a transformer, but it establishes the language-model training target: for a
sequence $x_0,\dots,x_{T-1}$, predict $x_1,\dots,x_T$.

\subsection*{Embedding and output head}
Input embeddings map token IDs to vectors. The output head maps vectors back to
vocabulary logits. Some language models tie these weights, using the same table
for input embeddings and output scoring. You do not need weight tying for the
first implementation, but the idea is useful: both sides of the model represent
tokens in the same learned vector space.

\subsection*{Exercise connection}
The exercise asks for the core lookup operation. The full project then places
that lookup inside a training loop with cross-entropy loss.

\medskip
\noindent\textit{Go deeper:} Karpathy's \texttt{makemore}
(\texttt{https://github.com/karpathy/makemore}) trains character-level language
models from token embeddings, the same baseline this module describes.
